\NSPage{109}%
\begin{EESegment}{EE-TGT-P109-001}
bound. In physical coordinates, the auxiliary averages are
\NSi{NS-F-P109-001}{F_2=Q^{-2A-1/2}\langle E_\theta\rangle_Y} and
\NSi{NS-F-P109-002}{F_1=Q^{-2A-1/2}\langle E_z\rangle_Y}, and these definitions agree from one
chart to another. For \NSi{NS-F-P109-003}{e\in\{1,2\}}, choose one fixed interior profile
\NSi{NS-F-P109-004}{\widehat b_e} with
\NSi{NS-F-P109-005}{\int x^e\widehat b_e(x)\,dx=1}. Set
\NSi{NS-F-P109-006}{b_e(r,z,t)=q^{-(e+1)/2}\widehat b_e(r/\sqrt q)}, and then set
\end{EESegment}
\NSFormula{NS-F-P109-007}{display}{none}
\[
 M_e=\int_0^\infty r^eF_e(r)\,dr,
 \qquad
 \sigma_e(r)=-r^{-e}\int_0^r(r')^e\bigl(F_e(r')-b_e(r')M_e\bigr)\,dr'.
\]
\begin{EESegment}{EE-TGT-P109-002}
These radial integrals are taken while \NSi{NS-F-P109-008}{(z,t)} stays fixed. The total weighted
moment of the integrand is zero because
\NSi{NS-F-P109-009}{\int r^e(F_e-b_eM_e)\,dr=M_e-M_e=0}. Starting the integral at zero fixes
the integration constant. Since the weighted moment is zero, \NSi{NS-F-P109-010}{\sigma_e}
becomes zero beyond the annulus that supports the source. Thus
\NSi{NS-F-P109-011}{\sigma_e} has compact support and
\NSi{NS-F-P109-012}{(\partial_r+e/r)\sigma_e=-F_e+b_eM_e}. Define the normalized tensor pair by
\NSi{NS-F-P109-013}{\Sigma=Q^{2A}(\sigma_2,\sigma_1)}. The weighted radial-integral estimate in
the proof of Proposition~\ref{prop:9.3} gives
\NSi{NS-F-P109-014}{\Sigma\in \mathsf M_{C_*-\kappa_s}}. This construction uses the auxiliary
averages \NSi{NS-F-P109-015}{F_e}, so the answer does not depend on any auxiliary torus
coordinate, even when the original fields have nonzero Fourier modes on that torus.
\end{EESegment}

\begin{EESegment}{EE-TGT-P109-003}
Apply the amplitude-correction operator \NSi{NS-F-P109-016}{\mathcal L^{\mathrm{as}}_{\mathrm{phys}}}
from \eqref{eq:7.35} to the pair \NSi{NS-F-P109-017}{\sigma=(\sigma_2,\sigma_1)}. This operator
adds up the localized operators in physical coordinates. Its normalized form
\NSi{NS-F-P109-018}{\mathcal L^{\mathrm{as}}\Sigma} lies in the supported class
\NSi{NS-F-P109-019}{\mathsf W_{B-\kappa_s}}. By \eqref{eq:7.36}, its prescribed symmetrized
cross covariance is \NSi{NS-F-P109-020}{\Sigma} with
\NSi{NS-F-P109-021}{w_0^{\tan}=W_0^{\mathrm{as}}}. Each local summand uses \eqref{eq:7.31} before the slow cutoff and
physical rescaling are applied. For each sign
\NSi{NS-F-P109-022}{\sigma\in\{+,-\}}, the formula divides its known numerator by the fixed,
positive primary amplitude \NSi{NS-F-P109-023}{2\sqrt{y_\sigma}}. It can therefore produce an
increment of either sign without taking a square root of the current covariance or of
\NSi{NS-F-P109-024}{\boldsymbol y+d\Sigma}. For the homogeneous amplitude equation, whose
source is zero, use the pressure coefficient. Build the velocity by taking the curl, as in
Lemma~\ref{lem:7.7}. The covariance identity pairs the two signs for one slow label. When their
covariances are added, each slow label is used once, together with the square of the slow
partition. The correction's quadratic self-interaction stays in the residual.
\end{EESegment}

\begin{EESegment}{EE-TGT-P109-004}
The nonzero harmonic errors now have these exponents.
\end{EESegment}
\NSFormula{NS-F-P109-025}{display}{none}
\[
\begin{array}{r|l}
\text{linear error} & B+1/2-4\kappa_s\\
\text{mean interaction} & B+0.4-\kappa_s\\
\text{cross with old exact waves} & B+1/2-2\kappa_s\\
\text{signed-correction self-interaction} & 2B-3\kappa_s.
\end{array}
\]
\begin{EESegment}{EE-TGT-P109-005}
For the auxiliary mean, split the old exact wave into its fixed primary transverse amplitude and a
remainder in \NSi{NS-F-P109-026}{\mathsf W_{0.68}}. More specifically, write the signed increment
as \NSi{NS-F-P109-027}{s=t_s+r_s}, where \NSi{NS-F-P109-028}{t_s} is the transverse amplitude
and \NSi{NS-F-P109-029}{r_s} is the curl remainder. Set
\NSi{NS-F-P109-030}{\mathscr B(a,b)=\langle\langle a\otimes b+b\otimes a\rangle_\theta\rangle_Y}.
Its exact change in covariance is
\end{EESegment}
\NSFormula{NS-F-P109-031}{display}{9.13}
\begin{equation}
\langle\Delta W\rangle_Y
=\mathscr B(w_0^{\tan},t_s)+\mathscr B(w-w_0^{\tan},t_s)
 +\mathscr B(w,r_s)+\langle\langle s\otimes s\rangle_\theta\rangle_Y.
\tag{9.13}\label{eq:9.13}
\end{equation}
\begin{EESegment}{EE-TGT-P109-006}
The radial--tangential components of the first term equal
\NSi{NS-F-P109-032}{\Sigma}, so their divergence cancels
\NSi{NS-F-P109-033}{F_e-b_eM_e}. Before taking the divergence, the tensor terms that remain have
these exponents.
\end{EESegment}
\NSFormula{NS-F-P109-034}{display}{none}
\[
\begin{array}{r|l}
\text{transverse correction with old-wave remainder}
  & B+0.68-\kappa_s=C_*+0.18-\kappa_s\\
\text{signed curl correction times old exact wave}
  & C_*+1/2-2\kappa_s\\
\text{signed-correction self-interaction}
  & 2B-2\kappa_s=C_*+\sigma_j-2\kappa_s.
\end{array}
\]
\begin{EESegment}{EE-TGT-P109-007}
These are exactly the remainder terms in the covariance expansion from
Corollary~\ref{cor:7.8}. Taking their radial divergence loses one power
\NSi{NS-F-P109-035}{\kappa_s}. Axial fluxes gain one when the tensor change has order at least
\NSi{NS-F-P109-036}{C_*-\kappa_s}. The changes in
\NSi{NS-F-P109-037}{g_r,p_m} have order at least
\NSi{NS-F-P109-038}{C_*-2\kappa_s}, and the effect of the pressure on
\NSi{NS-F-P109-039}{E_z} gains one. The moment-correction term
\NSi{NS-F-P109-040}{b_eM_e}, which has not been canceled, has order
\NSi{NS-F-P109-041}{C_*+1-\kappa_s} by \eqref{eq:9.12}.
\end{EESegment}
\NSPage{110}%
\begin{EESegment}{EE-TGT-P110-001}
Since \NSi{NS-F-P110-001}{0.18-2\kappa_s>0.17} and
\NSi{NS-F-P110-002}{\sigma_j-3\kappa_s>0.17}, the state after the pressure is rebuilt satisfies
\end{EESegment}
\NSFormula{NS-F-P110-003}{display}{9.14}
\begin{equation}
\langle E_\theta\rangle_Y,\ \langle E_z\rangle_Y\in \mathsf M_{C_*+0.17},
\qquad E_\theta,E_z\in \mathsf M_H,
\qquad (\mathcal P,\mathcal J_\theta,\mathcal J_z)\in \mathsf S_H,
\qquad H=C_*-2\kappa_s.
\tag{9.14}\label{eq:9.14}
\end{equation}

\begin{EESegment}{EE-TGT-P110-002}
\noindent\textbf{Step 3: Remove the nonconstant auxiliary means.} The auxiliary averages now have
the better order in \eqref{eq:9.14}. Their parts with zero auxiliary average still have order
\NSi{NS-F-P110-004}{H}, and the temporal mean inverse can cancel them. For the current exact
mean residuals, set
\NSi{NS-F-P110-005}{E_a^\circ=E_a-\langle E_a\rangle_Y}, with
\NSi{NS-F-P110-006}{a=\theta,z}. The temporal inverse from \eqref{eq:8.20} gives
\end{EESegment}
\NSFormula{NS-F-P110-007}{display}{none}
\[
\Delta v=-c_{i_0}^{-1}N_{i_0}^{-1}E_\theta^\circ,
\qquad
\gamma_d=-c_{i_0}^{-1}N_{i_0}^{-1}E_z^\circ,
\qquad
\Psi_*=\mathsf T_1\gamma_d,
\qquad
\Delta\beta=-\varepsilon\partial_Z\Psi_*,
\qquad
\Delta\gamma=(D_r+R^{-1})\Psi_*.
\]
\begin{EESegment}{EE-TGT-P110-003}
Here \NSi{NS-F-P110-008}{N_{i_0}^{-1}} is the inverse with zero average from
Lemma~\ref{lem:8.6}. The prescribed tangential increments lie in
\NSi{NS-F-P110-009}{\mathsf M_H}, and the radial velocity increment they produce lies in
\NSi{NS-F-P110-010}{\mathsf M_{H+1}}. Rebuild the pressure. The term
\NSi{NS-F-P110-011}{2V\Delta v/R} still has order \NSi{NS-F-P110-012}{H} in
\NSi{NS-F-P110-013}{\Delta g_r}, so \NSi{NS-F-P110-014}{\Delta p_m\in \mathsf M_H}. Every
other part of \NSi{NS-F-P110-015}{\Delta g_r} has one of the following orders. The fast-time
derivative of \NSi{NS-F-P110-016}{\Delta\beta} has order
\NSi{NS-F-P110-017}{H+1}. Each radial flux contains \NSi{NS-F-P110-018}{b}, an old radial
mean, or a new radial mean. The axial fluxes gain one. The other quadratic terms in the azimuthal
velocity have order at least \NSi{NS-F-P110-019}{H+0.9}. The viscous term contains a factor
\NSi{NS-F-P110-020}{\varepsilon} and loses at most two radial derivatives.
\end{EESegment}

\begin{EESegment}{EE-TGT-P110-004}
Write \NSi{NS-F-P110-021}{\Delta\gamma} for the axial increment that is actually produced, and
write \NSi{NS-F-P110-022}{\gamma_d} for the increment prescribed in \eqref{eq:8.20}. Since
\NSi{NS-F-P110-023}{\mathfrak t_*=-\varepsilon\partial_T+c_{i_0}N_{i_0}}, the fast-time
identities are
\end{EESegment}
\NSFormula{NS-F-P110-024}{display}{none}
\[
\begin{aligned}
c_{i_0}N_{i_0}\Delta v&=-E_\theta^\circ,\\
c_{i_0}N_{i_0}\Delta\gamma&=-E_z^\circ+\mathcal F_{\mathrm{ax}},\\
\mathcal F_{\mathrm{ax}}&=c_{i_0}N_{i_0}(\Delta\gamma-\gamma_d).
\end{aligned}
\]
\begin{EESegment}{EE-TGT-P110-005}
The last term is the flat remainder from reconstructing the axial increment in
Lemma~\ref{lem:8.6}. This mean update does not change the wave covariance. Subtracting the two
versions of \eqref{eq:8.3} gives
\end{EESegment}
\NSFormula{NS-F-P110-025}{display}{none}
\[
\begin{aligned}
E_{\theta,\mathrm{new}}-\langle E_\theta\rangle_Y
  ={}&-\varepsilon\partial_T\Delta v-\varepsilon(\Delta_0-R^{-2})\Delta v\\
     &+(D_r+2/R)\bigl((b+\beta)\Delta v+(V+v)\Delta\beta+\Delta\beta\Delta v\bigr)\\
     &+D_z\bigl((G+\gamma)\Delta v+(V+v)\Delta\gamma+\Delta\gamma\Delta v\bigr),\\
E_{z,\mathrm{new}}-\langle E_z\rangle_Y
  ={}&-\varepsilon\partial_T\Delta\gamma-\varepsilon\Delta_0\Delta\gamma+\mathcal F_{\mathrm{ax}}\\
     &+(D_r+1/R)\bigl((b+\beta)\Delta\gamma+(G+\gamma)\Delta\beta+\Delta\beta\Delta\gamma\bigr)\\
     &+D_z\bigl(2(G+\gamma)\Delta\gamma+(\Delta\gamma)^2+\Delta p_m\bigr).
\end{aligned}
\]
\begin{EESegment}{EE-TGT-P110-006}
Use \eqref{eq:9.9}, the bounds
\NSi{NS-F-P110-026}{|D^I b|\le C_I\varepsilon S_*^{b_I}} on the shell, and the orders of the
increments above. The four kinds of terms then have these lower bounds for their class exponents.
\end{EESegment}
\NSFormula{NS-F-P110-027}{display}{none}
\[
\begin{array}{c|cccc}
\text{term} & \text{slow time} & \text{radial flux} & \text{axial flux} & \text{viscosity}\\ \hline
\text{exponent} & H+1 & H+1-\kappa_s & H+1 & H+1-2\kappa_s.
\end{array}
\]
\begin{EESegment}{EE-TGT-P110-007}
After \NSi{NS-F-P110-028}{\mathcal F_{\mathrm{ax}}} is taken out of the axial component, the
displayed changes in the residual therefore lie in
\NSi{NS-F-P110-029}{\mathsf M_{H+1-2\kappa_s}}. Keep
\NSi{NS-F-P110-030}{\mathcal F_{\mathrm{ax}}} in the additive flat term. The change in the wave
lies in \NSi{NS-F-P110-031}{\mathsf W_H} because it interacts with the waves already added, whose
total obeys the \NSi{NS-F-P110-032}{\mathsf W_{1/2}} bound. The remainders from the modified,
compactly supported radial integral go into the additive flat term, while their actual velocities
continue to be used. The construction has now reached
\end{EESegment}
\NSFormula{NS-F-P110-033}{display}{none}
\[
E_\theta,E_z\in \mathsf M_{\min(C_*+0.17,H+1-2\kappa_s)},
\qquad
(\mathcal P,\mathcal J_\theta,\mathcal J_z)\in \mathsf S_H.
\]
\NSPage{111}%
\begin{EESegment}{EE-TGT-P111-001}
\noindent\textbf{Step 4: Enforce the three conditions needed for compact support.} The tangential
residuals have improved, but the compatibility defects still have order
\NSi{NS-F-P111-001}{H}. Apply the five-equation correction map \eqref{eq:8.25} to these defects
at order \NSi{NS-F-P111-002}{H}. The azimuthal increment and the prescribed axial increment are
slow. They are supported on the fixed support of the test-function profiles in the five-equation
correction system, and they lie in \NSi{NS-F-P111-003}{\mathsf M_H}. The radial velocity
increment they produce lies in \NSi{NS-F-P111-004}{\mathsf M_{H+1}}. The first two rows preserve
\eqref{eq:9.10}. The last three rows remove the linear contributions to the changes in
\NSi{NS-F-P111-005}{(\mathcal P,\mathcal J_\theta,\mathcal J_z)}. Rebuild the pressure from the
new fields that are actually used.
\end{EESegment}

\begin{EESegment}{EE-TGT-P111-002}
The earlier bounds for changes in the nonzero harmonics and the tangential mean residuals still
hold. Unlike a fast-time correction, a slow correction from this system has no fast-time term to
cancel. To estimate the defects that remain, subtract
\NSi{NS-F-P111-006}{2V\Delta v/R} from \NSi{NS-F-P111-007}{\Delta g_r}. What remains contains
the slow-time derivative of the radial velocity just produced, radial derivatives of
\NSi{NS-F-P111-008}{b\Delta\beta} and of products involving old or new radial means, axial flux
derivatives, products of \NSi{NS-F-P111-009}{\Delta v} with old or new means, and the radial
viscous term. Every one of these terms has exponent at least
\NSi{NS-F-P111-010}{H+0.9-2\kappa_s}. In
\NSi{NS-F-P111-011}{\Delta\mathcal J_\theta} and
\NSi{NS-F-P111-012}{\Delta\mathcal J_z}, the terms left after this system cancels the linear
ones are products of tangential means or moments of this same radial remainder. The exact formulas
are in \eqref{eq:8.27}. Hence
\end{EESegment}
\NSFormula{NS-F-P111-013}{display}{none}
\[
(\mathcal P,\mathcal J_\theta,\mathcal J_z)
 \in \mathsf S_{H+0.9-2\kappa_s}=\mathsf S_{C_*+0.9-4\kappa_s}.
\]
\begin{EESegment}{EE-TGT-P111-003}
These inequalities close the cycle.
\end{EESegment}
\NSFormula{NS-F-P111-014}{display}{none}
\[
\begin{gathered}
\min\{\tfrac12-3\kappa_s,\tfrac12-\kappa_s,B-\kappa_s,0.4\}\ge0.4,\\
\min\{\tfrac12-4\kappa_s,0.4-\kappa_s,\tfrac12-2\kappa_s,B-3\kappa_s\}
 \ge0.4-\kappa_s,\\
H-B=\tfrac12-2\kappa_s>0.1,
\qquad
\min\{0.17,1-4\kappa_s\}=0.17>0.1,\\
0.9-4\kappa_s>0.1.
\end{gathered}
\]
\begin{EESegment}{EE-TGT-P111-004}
Here \NSi{NS-F-P111-015}{B\ge0.7} was used. The orders of the wave, the full tangential mean,
and the defects all improve by \NSi{NS-F-P111-016}{0.1}. The smallest order of a signed
increment is \NSi{NS-F-P111-017}{B-\kappa_s\ge0.69999>0.68}, while the particular increments
start at \NSi{NS-F-P111-018}{B\ge0.7}. Every new increment in the mean velocity or mean pressure
has decay exponent greater than \NSi{NS-F-P111-019}{0.9}, and every increment in the radial mean
velocity has exponent greater than \NSi{NS-F-P111-020}{1.9}. A finite sum therefore still
satisfies \eqref{eq:9.9}. Reconstructing the pressure gives the claimed radial residual.
Proposition~\ref{prop:9.3} gives the source hypotheses for the next cycle, which completes the
induction.
\end{EESegment}
\end{proof}

\begin{EESegment}{EE-TGT-P111-005}
\subsection{One common domain and estimates for physical derivatives.}\label{sec:9.4}
\end{EESegment}
\begin{EESegment}{EE-TGT-P111-006}
To use Lemma~\ref{lem:5.4}, every finite state
\NSi{NS-F-P111-021}{(u^{[j]},p^{[j]})} must exist on the same domain. Lemma~\ref{lem:9.7}
provides that domain. The gain in decay must also survive differentiation in the physical
variables. More exactly, once the derivative order is fixed, the loss in powers of
\NSi{NS-F-P111-022}{q} must not depend on the correction stage. The constants in the bounds may
still depend on that stage, as Lemma~\ref{lem:9.8} allows.
\end{EESegment}

\begin{lemma}\label{lem:9.7}
\begin{EESegment}{EE-TGT-P111-007}
There is a number \NSi{NS-F-P111-023}{q_{\mathrm{big}}>0} that does not depend on the correction
stage or the derivative order, and every finite partial sum above, before its summation cutoffs are
applied, is defined on \NSi{NS-F-P111-024}{0<q<q_{\mathrm{big}}}. Once the stage and the order of
an output amplitude derivative are fixed, constructing and estimating that derivative needs only
finitely many derivatives of the input amplitudes. The constants may depend on the correction
stage.
\end{EESegment}
\end{lemma}

\begin{proof}
\begin{EESegment}{EE-TGT-P111-008}
Choose \NSi{NS-F-P111-025}{0<q_{\mathrm{big}}\le q_*} once and for all so that the pulse estimates
hold whenever \NSi{NS-F-P111-026}{q<q_{\mathrm{big}}}. Make this choice with the fixed
background, the stated positive lower bound for \NSi{NS-F-P111-027}{|n_\Phi|}, the primary
covariance, and the fixed supports of the normalized mean-correction profiles. Every later
inhomogeneous amplitude problem is a linear equation whose coefficients are these fixed ones and
whose source is already known. The same threshold works for every harmonic because
\eqref{eq:7.18} writes its propagator as the fundamental propagator times a forward damping
factor whose magnitude is at most one. At each
\NSPage{112}%
fixed derivative order, differentiation puts powers of the harmonic integer into the constants and
allows polynomial growth in \NSi{NS-F-P112-001}{S_*}. It does not call for a smaller
\NSi{NS-F-P112-002}{q}-threshold, as Corollary~\ref{cor:7.3} shows.
\end{EESegment}

\begin{EESegment}{EE-TGT-P112-001}
The signed updates divide by the original \NSi{NS-F-P112-003}{2\sqrt{y_\sigma}} in
\eqref{eq:7.31}, and the finite-dimensional correction system uses the fixed linear map from
Lemma~\ref{lem:8.7}. Pressure, the temporal inverses, and the modified radial integrals that define
the mean potentials are all fixed linear maps. The primary construction fixes every inverse
coefficient and every denominator. These finite operations are therefore defined on the same
domain for sources of any size, including sources near
\NSi{NS-F-P112-004}{q_{\mathrm{big}}}. Each cycle improves the decay exponent on that whole
domain.
\end{EESegment}

\begin{EESegment}{EE-TGT-P112-002}
The number of input derivatives needed at a given stage can be counted with a directed acyclic
graph. Each vertex is an intermediate expression, and each arrow records one of the sums,
products, derivatives, or proved inverses used to build a later expression from an earlier one.
Directed means that these arrows point from input toward output. Acyclic means that no chain of
arrows loops back to where it started. The source vertices, where no earlier expression feeds in,
are the input fields. For a source vertex \NSi{NS-F-P112-005}{a}, let
\NSi{NS-F-P112-006}{D_{v,a}(m)} be any sufficient number of derivatives of the input
\NSi{NS-F-P112-007}{a} for bounding, through derivative order
\NSi{NS-F-P112-008}{m}, the expression at vertex \NSi{NS-F-P112-009}{v}. Suppose the operation
at \NSi{NS-F-P112-010}{v} needs \NSi{NS-F-P112-011}{d_{v,w}} extra derivatives of its input at
\NSi{NS-F-P112-012}{w}. Define
\end{EESegment}
\NSFormula{NS-F-P112-013}{display}{none}
\[
D_{v,a}(m)=\max_{w\to v}D_{w,a}(m+d_{v,w}),
\]
\begin{EESegment}{EE-TGT-P112-003}
starting with \NSi{NS-F-P112-014}{D_{b,a}(m)=m} when the source vertices satisfy
\NSi{NS-F-P112-015}{b=a}, and with \NSi{NS-F-P112-016}{D_{b,a}(m)=0} when
\NSi{NS-F-P112-017}{b\ne a}. At every fixed correction stage the graph has only finitely many
vertices, so every derivative count produced this way is finite. Differentiating a pulse equation
to order \NSi{NS-F-P112-018}{m} gives an equation containing amplitude derivatives of order at
most \NSi{NS-F-P112-019}{m}; the lower-order derivatives are already controlled. The torus
inverses need finitely many more amplitude derivatives, while the modified radial integrals keep
the stated derivative bounds. To make the cutoff remainder flat to a chosen power, integrate the
same exact integral by parts only finitely many times. These counts tell exactly how many source
derivatives need estimates. The explicit \NSi{NS-F-P112-020}{\varepsilon} losses instead lower
the decay exponent. For example, estimates for more amplitude derivatives of
\NSi{NS-F-P112-021}{N^{-1}D_rf} still have the one explicit radial loss belonging to
\NSi{NS-F-P112-022}{D_r}. First construct and estimate every finite-stage coefficient as a
function of independent slow and auxiliary variables. Then set the auxiliary variables equal to
the phase map \NSi{NS-F-P112-023}{Y=Y(r,t)} and differentiate with respect to
\NSi{NS-F-P112-024}{(x,t)}.
\end{EESegment}
\end{proof}

\begin{EESegment}{EE-TGT-P112-004}
Put the base and the finite initialization together as
\NSi{NS-F-P112-025}{U_0=(u^{[0]},p^{[0]})}. For \NSi{NS-F-P112-026}{j\ge1}, let
\NSi{NS-F-P112-027}{A_j^w} be the sum of the physical wave-potential increments made during the
cycle from stage \NSi{NS-F-P112-028}{j-1} to stage \NSi{NS-F-P112-029}{j}. Let
\NSi{NS-F-P112-030}{\Psi_j} be the sum of that cycle's physical azimuthal-potential coefficients,
and let \NSi{NS-F-P112-031}{b_j} be the sum of its direct azimuthal velocity coefficients in
physical variables. In the rest of the summation argument,
\NSi{NS-F-P112-032}{B_j} means the azimuthal vector from \eqref{eq:5.29}. Define
\end{EESegment}
\NSFormula{NS-F-P112-033}{display}{9.15}
\begin{equation}
\begin{gathered}
A_j=A_j^w+\Psi_je_\theta,
\qquad B_j=b_je_\theta,\\
p_j=p^{[j]}-p^{[j-1]},
\qquad Z_j=(A_j,B_j,p_j),\\
\Delta u_j=\curl A_j+B_j=u^{[j]}-u^{[j-1]},\\
(u^{[J]},p^{[J]})=U_0+\sum_{j=1}^{J}(\Delta u_j,p_j).
\end{gathered}
\tag{9.15}\label{eq:9.15}
\end{equation}
\NSPage{113}%
\begin{EESegment}{EE-TGT-P113-001}
Thus \NSi{NS-F-P113-001}{Z_j} is the tuple from \eqref{eq:5.29}, without the optional stress
entry. In a fixed \NSi{NS-F-P113-002}{Q} chart, the linearity of
\NSi{NS-F-P113-003}{\mathsf T_0} and the fixed choice of \NSi{NS-F-P113-004}{\rho} give this
explicit formula for the pressure increment.
\end{EESegment}
\NSFormula{NS-F-P113-005}{display}{9.16}
\begin{equation}
Q^{2A}p_j=p_w^{[j]}-p_w^{[j-1]}
 +\mathsf T_0\bigl(g_r^{[j]}-g_r^{[j-1]}-\rho(\mathcal P^{[j]}-\mathcal P^{[j-1]})\bigr).
\tag{9.16}\label{eq:9.16}
\end{equation}
\begin{EESegment}{EE-TGT-P113-002}
Here \NSi{NS-F-P113-006}{p_w^{[j]},g_r^{[j]},\mathcal P^{[j]}} are the normalized quantities
from \eqref{eq:9.7} and \eqref{eq:9.6}. The pressure differences from all the steps in between
telescope, leaving \eqref{eq:9.16}.
\end{EESegment}

\begin{lemma}\label{lem:9.8}
\begin{EESegment}{EE-TGT-P113-003}
For every order \NSi{NS-F-P113-007}{m} of Cartesian space--time derivatives,
\end{EESegment}
\NSFormula{NS-F-P113-008}{display}{9.17}
\begin{equation}
|Z_j|_m+|\Delta u_j|_m
 \le C_{j,m}q^{g_j-\ell_m}(1+|\log q|)^{P_{j,m}},
\qquad
g_j=\frac h{10}j\longrightarrow\infty,
\tag{9.17}\label{eq:9.17}
\end{equation}
\begin{EESegment}{EE-TGT-P113-004}
where \NSi{NS-F-P113-009}{\ell_m} does not depend on \NSi{NS-F-P113-010}{j} and includes the
extra spatial derivative used to recover velocities from the wave and mean vector potentials. The
finite partial sums satisfy
\end{EESegment}
\NSFormula{NS-F-P113-011}{display}{none}
\[
|(u^{[j]},p^{[j]})|_m
 \le C_{j,m}q^{-K_m}(1+|\log q|)^{P_{j,m}},
\]
\begin{EESegment}{EE-TGT-P113-005}
with \NSi{NS-F-P113-012}{K_m} independent of \NSi{NS-F-P113-013}{j}. Their residuals satisfy
\end{EESegment}
\NSFormula{NS-F-P113-014}{display}{9.18}
\begin{equation}
|\mathscr R(u^{[j]},p^{[j]})|_m
 \le C_{j,m}q^{h\sigma_j-K_m}(1+|\log q|)^{P_{j,m}}+E_{j,m},
\qquad
E_{j,m}\le C_{j,m,N}q^N\quad\text{for every }N,
\tag{9.18}\label{eq:9.18}
\end{equation}
\begin{EESegment}{EE-TGT-P113-006}
and here too \NSi{NS-F-P113-015}{K_m} is independent of \NSi{NS-F-P113-016}{j}.
\end{EESegment}
\end{lemma}

\begin{proof}
\begin{EESegment}{EE-TGT-P113-007}
On the supports of the perturbations, \NSi{NS-F-P113-017}{r\asymp\sqrt Q} and
\NSi{NS-F-P113-018}{q\asymp Q}. While differentiating in one fixed chart, keep
\NSi{NS-F-P113-019}{Q} fixed. For an amplitude coefficient, the graph operators in
\eqref{eq:6.6} lose the following powers of \NSi{NS-F-P113-020}{Q}.
\end{EESegment}
\NSFormula{NS-F-P113-021}{display}{9.19}
\begin{equation}
\begin{array}{r|l}
\text{radial physical derivative} & 1/2+h\kappa_s\\
\text{axial physical derivative} & D=1/2-h\\
\text{time physical derivative} & 1+h\\
\text{derivative of the cylindrical frame} & 1/2.
\end{array}
\tag{9.19}\label{eq:9.19}
\end{equation}
\begin{EESegment}{EE-TGT-P113-008}
For example, the coefficient of an auxiliary-coordinate derivative in the physical radial
derivative is
\NSi{NS-F-P113-022}{r^{d_r-1}\Lambda_g^{i_0}=O(Q^{-1/2}\varepsilon^{-\kappa_s}S_*^P)},
while \NSi{NS-F-P113-023}{T_g^{i_0}=O(Q^{-1-h}S_*^P)}. Each additional fixed derivative lowers
the power of \NSi{NS-F-P113-024}{Q} by a fixed amount. The bounds for these exponents stay
uniform when the chart indices differ by only a bounded amount.
\end{EESegment}

\begin{EESegment}{EE-TGT-P113-009}
The label phase satisfies
\NSi{NS-F-P113-025}{|\nabla_x\Phi|\le CQ^{-1/2}S_*^P} and
\NSi{NS-F-P113-026}{|\partial_t\Phi|\le CQ^{-1-h}S_*^P}. The axial term
\NSi{NS-F-P113-027}{p_zZ/\varepsilon} costs
\NSi{NS-F-P113-028}{Q^{-D}\varepsilon^{-1}=Q^{-1/2}}. More generally, for every fixed pair of
nonnegative integers \NSi{NS-F-P113-029}{a,b} with \NSi{NS-F-P113-030}{a+b\ge1},
\eqref{eq:7.3} and \eqref{eq:6.12} give
\end{EESegment}
\NSFormula{NS-F-P113-031}{display}{none}
\[
|\nabla_x^a\partial_t^b\Phi_\gamma|
 \le C_{a,b}Q^{-a/2-b(1+h)}S_*^{P_{a,b}},
\]
\begin{EESegment}{EE-TGT-P113-010}
uniformly over all bands and labels. In each chart, the pulse coordinate has no physical spatial
derivatives and has one fixed first physical time derivative. Every other factor has the slow
derivative bounds already stated. Set
\end{EESegment}
\NSFormula{NS-F-P113-032}{display}{none}
\[
s_x=\tfrac12+\tfrac h2,
\qquad
s_t=1+\tfrac{3h}{2}.
\]
\begin{EESegment}{EE-TGT-P113-011}
Since \NSi{NS-F-P113-033}{k\le2Q^{-h/2}}, the chain rule gives
\end{EESegment}
\NSFormula{NS-F-P113-034}{display}{none}
\[
|\nabla_x^a\partial_t^b e^{ikm\Phi}|
 \le C_{a,b,m}Q^{-as_x-bs_t}S_*^{P_{a,b,m}}.
\]
\begin{EESegment}{EE-TGT-P113-012}
These costs are at least as large as the ones in \eqref{eq:9.19}, so they control them. At any fixed
stage there are only finitely many harmonic integers, and they change only the constants. A
quadratic residual can at most double the current harmonic range, while inverse and curl operations
keep each integer unchanged. The harmonic range is therefore finite independently of the band.
\end{EESegment}
\NSPage{114}%
\begin{EESegment}{EE-TGT-P114-001}
For a class exponent \NSi{NS-F-P114-001}{\alpha}, Leibniz' rule now bounds the Cartesian
space--time derivatives of the normalized coefficient by
\NSi{NS-F-P114-002}{CQ^{h\alpha-as_x-bs_t}S_*^P}. The factors
\NSi{NS-F-P114-003}{\sqrt\zeta\,\delta^{-M}} and
\NSi{NS-F-P114-004}{\zeta\delta^{-M}} are bounded on the closed shell for every fixed
\NSi{NS-F-P114-005}{M}, so they do not cause any further loss in powers. The rescaling factors
for physical velocity, pressure, wave potential, and mean potential are, respectively,
\NSi{NS-F-P114-006}{Q^{-A},Q^{-2A},Q^{1/2-A},Q^{1/2-A}}. Each of them is bounded by
\NSi{NS-F-P114-007}{Q^{-2A}}. The normalized chart wave potential already contains the factor
\NSi{NS-F-P114-008}{(km)^{-1}}, before physical rescaling. Recovering a velocity from a wave or
mean vector potential costs one more spatial derivative. One common choice that is large enough is
\end{EESegment}
\NSFormula{NS-F-P114-009}{display}{none}
\[
\ell_m=2A+(m+1)(1+3h/2),
\]
\begin{EESegment}{EE-TGT-P114-002}
Every increment kept at stage \NSi{NS-F-P114-010}{j} has normalized exponent at least
\NSi{NS-F-P114-011}{j/10}, and its physical prefactor is bounded by
\NSi{NS-F-P114-012}{Q^{-2A}}. This proves \eqref{eq:9.17}.
\end{EESegment}

\begin{EESegment}{EE-TGT-P114-003}
Apply the same calculation to \eqref{eq:9.9}, to the fixed background, and to \eqref{eq:9.8}.
Converting the residual from normalized coordinates to physical coordinates adds one fixed power of
\NSi{NS-F-P114-013}{q}. The rest of the radial residual consists of
\NSi{NS-F-P114-014}{-\rho\mathcal P}, written in physical coordinates, and its cutoff remainder.
Equation~\eqref{eq:9.5} gives \NSi{NS-F-P114-015}{E_{j,m}} uniformly over all labels at each
fixed stage. This proves \eqref{eq:9.18} and the bound for the whole correction. Where two
coordinate representations overlap, the fields agree after the prescribed changes of torus
coordinates and relabeling. Their smooth zero extensions then carry the estimates across support
and band boundaries.
\end{EESegment}
\end{proof}

\begin{EESegment}{EE-TGT-P114-004}
\subsection{Building the local field.}\label{sec:9.5}
\end{EESegment}
\begin{EESegment}{EE-TGT-P114-005}
The finite partial sums \NSi{NS-F-P114-016}{(u^{[j]},p^{[j]})} now have residuals whose decay can
be made arbitrarily high on one common domain. Lemmas~\ref{lem:9.7} and~\ref{lem:9.8} also give
the bounds for physical derivatives that Lemma~\ref{lem:5.4} needs. This subsection sums the
potentials with cutoffs that shrink from one stage to the next, producing the local velocity
\NSi{NS-F-P114-017}{u_{\mathrm{loc}}} from Proposition~\ref{prop:9.9}. The remaining work is to
check that this sum keeps the endpoint regularity, the exterior heat flow, and the inner growth in
Theorem~\ref{thm:3.1}(ii)--(iv).
\end{EESegment}

\begin{proposition}\label{prop:9.9}
\begin{EESegment}{EE-TGT-P114-006}
There are smooth fields \NSi{NS-F-P114-018}{(u_{\mathrm{loc}},p_{\mathrm{loc}})} on
\NSi{NS-F-P114-019}{\Omega_*}, obtained by summing the finite corrections above, such that
\NSi{NS-F-P114-020}{\divergence u_{\mathrm{loc}}=0} and every conclusion of
Theorem~\ref{thm:3.1} holds.
\end{EESegment}
\end{proposition}

\begin{proof}
\begin{EESegment}{EE-TGT-P114-007}
Apply the summation lemma with
\NSi{NS-F-P114-021}{\mathcal F(u,p)=\mathscr R(u,p)}. Use
\NSi{NS-F-P114-022}{q_*=q_{\mathrm{big}}} from Lemma~\ref{lem:9.7}, shrinking the earlier local
thresholds to this common value. The inputs are the velocity field and the pressure field. Step 1
checks the hypotheses of Lemma~\ref{lem:5.4}. Steps 2--5 prove
Theorem~\ref{thm:3.1}(i)--(iv), in that order. They prove Cartesian regularity, endpoint
regularity, the exterior formula together with the flat residual, and the inner swirl asymptotic.
\end{EESegment}

\begin{EESegment}{EE-TGT-P114-008}
\noindent\textbf{Step 1: Choose representatives and check the summation hypotheses.} Take
\NSi{NS-F-P114-023}{U_0} to be the fixed slow base together with the finite initialization block.
If needed, apply one cutoff to that block inside \NSi{NS-F-P114-024}{q<q_{\mathrm{big}}}. Apply it
to the block's wave and mean potentials before differentiation, and also to its direct angular means
and pressures. The cutoff equals one near \NSi{NS-F-P114-025}{q=0}, so every finite state still
agrees there with its original expression.
\end{EESegment}

\begin{EESegment}{EE-TGT-P114-009}
For \NSi{NS-F-P114-026}{j\ge1}, use the representatives
\NSi{NS-F-P114-027}{Z_j=(A_j,B_j,p_j)} from \eqref{eq:9.15}. Their mean-potential part gives
\end{EESegment}
\NSFormula{NS-F-P114-028}{display}{none}
\[
\curl(\Psi_je_\theta)=(-\partial_z\Psi_j,0,(\partial_r+r^{-1})\Psi_j).
\]
\begin{EESegment}{EE-TGT-P114-010}
Since \NSi{NS-F-P114-029}{B_j=b_je_\theta} and
\NSi{NS-F-P114-030}{\partial_\theta b_j=0}, multiplying \NSi{NS-F-P114-031}{B_j} by any
function of \NSi{NS-F-P114-032}{q(z,t)} keeps its divergence equal to zero. Every one of these
representatives extends smoothly by zero and vanishes near the axis for each fixed
\NSi{NS-F-P114-033}{t<1}. Because it vanishes near the axis, its Cartesian extension is smooth
there.
\end{EESegment}

\begin{EESegment}{EE-TGT-P114-011}
Lemma~\ref{lem:9.7} gives one domain for every finite partial sum before the summation cutoffs are
applied. The similarity identities give \eqref{eq:5.28}, and
\NSi{NS-F-P114-034}{q\ge1-t} gives the local lower bound that the summation lemma needs.
Lemma~\ref{lem:9.8} gives \eqref{eq:5.30} with
\NSi{NS-F-P114-035}{g_j=hj/10}, and the losses in the exponents do not depend on
\NSi{NS-F-P114-036}{j}.
\end{EESegment}
